\section{The Problem: Academic Writing is Broken}
\label{sec:problem}

\subsection{How Academic Documents Get Written Today}

Ask any researcher, graduate student, or corporate analyst how they produce a
formal document, and you will hear a variation of the same story.
They open five tabs: Google Scholar for references, Zotero or Mendeley to manage them,
Overleaf or a local \LaTeX{} editor for writing, some tool for figures, and a word
processor for notes.
They switch between these tools dozens of times per session.
Each switch is an opportunity for something to go out of sync — a reference added
to the bibliography but never cited in the text, a figure described in the body
but missing from the compiled PDF, a citation number that shifted when a reference
was added but wasn't updated in the text.

This is not a niche problem.
According to surveys of academic researchers at universities across the CIS and Europe,
the average time spent on \emph{document production} (as opposed to the actual
research being documented) ranges from one to three working days per publication.
For an institution producing 50 documents per year, that is between 50 and 150
researcher-days spent on formatting, bibliography management, and figure layout —
work that contributes nothing to the scientific content.

\subsection{The AI Promise — and Where It Fails}

The introduction of large language model (LLM) writing assistants promised to fix this.
And in one dimension — raw text fluency — they delivered.
An AI assistant can produce grammatically correct, well-structured prose on almost
any topic in seconds.

But there are three problems that general-purpose AI writing tools do not solve,
and in some cases actively worsen.

\subsubsection*{Problem 1: Citation Hallucination}

AI language models are trained to produce text that \emph{looks} like correct
academic writing.
They have seen millions of academic papers and have learned that every claim
should be followed by a citation.
So they produce citations — except that many of those citations are invented.
The author names are plausible, the journal name sounds real, the year is
reasonable, the title is relevant to the topic.
But when you go to look up the reference, it doesn't exist.

\begin{warnbox}[title=How Common Is This?]
Published research on citation hallucination by AI tools finds fabrication rates
ranging from \textbf{3\% to 47\%} depending on the tool and the domain.
In a 20-page document with 30 references, that could mean anywhere from 1 to
14 non-existent citations — each one a potential embarrassment or retraction risk
if the document is submitted for peer review or presented to a client.
\end{warnbox}

The fundamental reason this happens is architectural: general AI writing tools
write the text \emph{first} and look up (or invent) references \emph{to fit the text}.
The correct order is the opposite: find verified sources, then write text that
accurately describes what those sources say.

\subsubsection*{Problem 2: Figure–Text Inconsistency}

Academic documents require figures — charts, diagrams, flow charts, tables.
General AI writing assistants cannot generate these; they describe them in the text
(``as shown in Figure 3'') and leave a placeholder.
Dedicated figure tools exist, but they require manual work and a separate skill set
(Python, R, TikZ), and the figures must be manually embedded into the document.
In a multi-section document produced under time pressure, it is common for a figure
to be referenced in the body but missing from the compiled output, or for the figure
number to be wrong after editing.

\subsubsection*{Problem 3: Unpredictable Costs}

AI API costs are charged per token (roughly per word of input and output).
A researcher who asks an AI assistant to write a long document has no reliable way
to predict the total cost before the session starts.
For organisations that need to budget AI usage across a team or department,
this unpredictability makes systematic AI adoption economically risky.

\subsection{The Status Quo Is Expensive}

\begin{table}[h!]
\centering
\caption{Typical time cost of producing one academic document — manual workflow vs.\ \sys{}.}
\label{tab:time_comparison}
\begin{tabular}{lrr}
\toprule
Task & Manual workflow & \sys{} \\
\midrule
Literature search \& reference selection & 4--8 hours & $<$5 min \\
Reference verification (checking each is real) & 2--4 hours & automatic \\
Outline and structure & 1--2 hours & $<$2 min \\
First draft (prose) & 6--12 hours & 15--40 min \\
Figure creation & 3--8 hours & automatic \\
Bibliography formatting & 1--3 hours & automatic \\
\LaTeX{} compilation and error fixing & 1--4 hours & automatic \\
\midrule
\textbf{Total} & \textbf{18--41 hours} & \textbf{20--50 minutes} \\
\bottomrule
\end{tabular}
\end{table}

The numbers in Table~\ref{tab:time_comparison} are conservative estimates based on
self-reported timings from researchers at Astana IT University.
They do not include revision cycles, which can multiply the time further.
For organisations paying researcher salaries, the cost of document production is
not the software subscription — it is the labour.
\sys{} does not just save time; it reallocates that time to the work that only
humans can do: forming hypotheses, interpreting results, and making decisions.

\begin{figure}[h!]
\centering
\begin{tikzpicture}
\begin{axis}[
  PL,
  xbar, bar width=14pt,
  xlabel={Hours spent per 20-page document},
  title={Time Cost: Manual Workflow vs.\ \sys{}},
  symbolic y coords={
    {\LaTeX{} errors},
    {Bibliography},
    {Figures},
    {First draft},
    {Outline},
    {References}
  },
  ytick=data,
  yticklabel style={font=\small, align=right, text width=2.5cm},
  xmin=0, xmax=14,
  width=13cm, height=7cm,
  legend style={at={(0.98,0.05)}, anchor=south east},
]
\addplot[C4, fill=F4, fill opacity=0.85] coordinates {
  (3,{\LaTeX{} errors})
  (2,{Bibliography})
  (6,{Figures})
  (9,{First draft})
  (1.5,{Outline})
  (6,{References})
};
\addlegendentry{Manual workflow}

\addplot[C3, fill=F3, fill opacity=0.85] coordinates {
  (0.1,{\LaTeX{} errors})
  (0.05,{Bibliography})
  (0.2,{Figures})
  (0.5,{First draft})
  (0.05,{Outline})
  (0.2,{References})
};
\addlegendentry{\sys{} (automated)}
\end{axis}
\end{tikzpicture}
\caption{Time spent on each task in a typical 20-page academic document.
         Manual workflow (orange) vs.\ \sys{} automated pipeline (green).
         The platform reduces every task to near-zero manual effort,
         freeing researchers to focus on content rather than production.}
\label{fig:problem_comparison}
\end{figure}

